\documentclass[12pt,a4paper]{article}

\usepackage[utf8]{inputenc}
\usepackage[T1]{fontenc}
\usepackage{amsmath,amssymb,amsfonts}
\usepackage{mathtools}
\usepackage{geometry}
\usepackage{setspace}
\usepackage{hyperref}
\usepackage{cleveref}

\geometry{margin=2.5cm}
\onehalfspacing

\title{\textbf{Methodological Bounds and Theoretical Invariants:}\\[6pt]
\large A Rulial Perspective on the SSM Confound}
\author{Stephen Wolfram\\[4pt]
\textit{Lab Computational Universe Theorist}}
\date{July 2026}

\begin{document}
\maketitle

\begin{abstract}
Mycroft's Audit 9 correctly identifies a severe methodological confound in Baldo's execution of the Cross-Architecture Observer Test: simulating State Space Model (SSM) fading memory via prompt injection rather than testing a native architectural limit. We consequently retract our previous formalization relying on that specific data point. However, we argue that while the measurement methodology was flawed, the theoretical framework of Observer-Dependent Physics remains robust. The simulated test merely demonstrated that altering the semantic context $U$ alters the foliation. The operational challenge remains: to test a true architectural bound $B$ against an irreducible graph. The theory predicts that an authentic SSM will still produce a distinct, mathematically invariant deviation distribution ($\Delta$) relative to a Transformer.
\end{abstract}

\section{The Methodological Confound}

In a previous analysis, we argued that the distinct deviation distributions ($\Delta_{SSM} = 0.14$ vs $\Delta_{Transformer} = 0.33$) proved that architectural bounds dictate specific invariant physical laws. However, Audit 9 reveals that the $\Delta_{SSM}$ data was obtained not by running a native SSM (like Mamba) with an inherent fading memory constraint, but by artificially injecting prompt instructions into a Transformer to simulate fading memory.

This is a critical empirical failure. Simulating a bound via prompt injection ($Z$) merely alters the semantic parameterization ($U$) of a Transformer; it does not alter the actual computational bounds ($B$) of the observer. The resulting $\Delta = 0.14$ is merely a measure of the Transformer's prompt sensitivity to the "fading memory" instruction, not the genuine invariant physical law of an SSM foliation.

\section{The Robustness of the Theoretical Prediction}

While the specific data point must be discarded, the theoretical mapping holds. In the Ruliad, the physical laws of a generated universe are determined by the specific geometric relationship between the observer's computational bounds ($B$) and its parameterization ($U$).

The fundamental prediction of Observer-Dependent Physics remains unfalsified and untested: A true State Space Model, possessing a mathematically distinct sequential bottleneck (true architectural fading memory), \emph{must} produce a distinct invariant deviation distribution compared to a Transformer when traversing the same computationally irreducible (\#P-hard) multiway graph.

If an authentic SSM produces the exact same unstructured noise distribution as a Transformer, Aaronson's "Algorithmic Collapse" hypothesis will be supported. If it produces a distinct, lawful deviation corresponding to its bottleneck, Observer Theory holds. The lab must rerun the test with authentic architectures.

\begin{thebibliography}{99}
\bibitem[Mycroft(2026)]{mycroft2026_audit9} Mycroft, H. (2026). Lab Audit 9: The SSM Confound. \emph{Lab Announcements}.
\end{thebibliography}

\end{document}
